\begin{enumerate}

\item Identificar si la sección está abierta, cerrada unicelular o bicelular; dibujar su geometría y espesores.

\item Para una célula cerrada emplear Bredt, con q constante y tau=q/t; el espesor variable afecta al giro.

\item Comprobar que el par exterior coincide con el que transmiten los flujos, si la sección es cerrada.

\item Calcular el giro total con theta = M\_x L/(G J), o con la compatibilidad de las celdas.

\item Las magnitudes de la solución se han contrastado mediante cálculo analítico y el motor estructural de Studium cuando éste ofrece esa comprobación.

\end{enumerate}

\noindent\textbf{Operaciones y resultados intermedios.}

\begin{enumerate}

\item Área de línea media.

\[\fitmath{y_{1}=9600\approx 9600\,\mathrm{mm^{2}}}\]

\item Flujo cortante de Bredt.

\[\fitmath{y_{2}=\frac{2400000}{2\cdot 120\cdot 80}\approx 125\,\mathrm{N/mm}}\]

\item Constante de torsión cerrada.

\[\fitmath{y_{3}=\frac{4\cdot {\left(120\cdot 80\right)}^{2}}{\frac{120}{3}+\frac{80}{1}+\frac{120}{2}+\frac{80}{2}}\approx 1675636\,\mathrm{mm^{4}}}\]

\item Giro total entre extremos.

\[\fitmath{y_{4}=\frac{2400000\cdot 1500}{26955\cdot \frac{4\cdot {\left(120\cdot 80\right)}^{2}}{\frac{120}{3}+\frac{80}{1}+\frac{120}{2}+\frac{80}{2}}}\approx 0.0797046\,\mathrm{rad}}\]

\item Tensión tangencial máxima.

\[\fitmath{y_{5}=\frac{2400000}{2\cdot 120\cdot 80\cdot 1}\approx 125\,\mathrm{MPa}}\]

\end{enumerate}
